\section{\termAbstractFloerComplex{}及其同调}
\label{sec:abstract-floer-complex-homology}

本节给出所有已知的一般\personFloer{}同调构造背后的纯代数框架。我们沿用
\stfhIndex{I190}{Eshe}{厄舍（Usher）}
M. \personUsher{} (Ush08) 的叙述，并略作扩充。

\begin{definition}
\label{def:graded-filtered-floer-novikov-complex}
环 $R$ 上的一个\termGradedFilteredFloerNovikovComplex{} $\mathfrak c$ 由如下数据组成。
\begin{enumerate}[label=(\arabic*)]
  \item 一个群 $\Gamma$ 和一个（带离散拓扑的）主 $\Gamma$-丛
  \[
    P\longrightarrow S,
  \]
  其中
  \begin{enumerate}[label=(\alph*)]
    \item $S$ 是有限集；
    \item $\Gamma$ 是以乘法记号书写的有限生成阿贝尔群，而 $P$ 是一个
    $\Gamma$-集（即带 $\Gamma$-作用的集合）。
  \end{enumerate}
  \item 一个作用量泛函
  \[
    \mathcal A:P\longrightarrow\mathbb R
  \]
  和一个\termPeriodHomomorphism{}
  \footnote{译者注：原书把 $\Gamma$ 称为“群胚”及“有限生成阿贝尔群胚”，并把
  $\omega$ 称为“周期群胚同态”；但本定义使用全局乘法与群作用，本节后文还使用
  $\ker\omega$、群环和群的正合序列，均要求 $\Gamma$ 是群。原书又把 $P$ 写成
  未定义的 $G$-集（或 $G$-作用）。译文依这些代数结构统一为群 $\Gamma$、
  $\Gamma$-集和群同态。}
  \[
    \omega:\Gamma\longrightarrow\mathbb R,
  \]
  满足
  \[
    \mathcal A(g\cdot p)=\mathcal A(p)-\omega(g),
    \qquad g\in\Gamma,\ p\in P.
  \]

  \item 一个分次
  \[
    \operatorname{gr}:P\longrightarrow\mathbb Z
  \]
  和一个\termDegreeHomomorphism{}
  \[
    d:\Gamma\longrightarrow\mathbb Z,
  \]
  满足
  \footnote{译者注：原书在下式写成加号；但第 13.9 节由\personMaslov{}--\personMorse{}指标直接
  算得 $\operatorname{gr}([\widehat p,u\#w])
  =\operatorname{gr}([\widehat p,w])-\mu(u)$，并随即把
  $\operatorname{gr}(g\cdot p)=\operatorname{gr}(p)-d(g)$ 称为抽象定义所要求的
  公式。译文依这个具体构造及同章后文的明确复述补正为减号。}
  \[
    \operatorname{gr}(g\cdot p)
    =\operatorname{gr}(p)-d(g),
    \qquad g\in\Gamma,\ p\in P.
  \]

  \item 一个映射
  \[
    n:P\times P\longrightarrow R,
  \]
  满足如下条件。
  \begin{enumerate}[label=(\alph*)]
    \item 除非同时有
    \[
      \mathcal A(p)>\mathcal A(p')
      \quad\text{和}\quad
      \operatorname{gr}(p)=\operatorname{gr}(p')+1,
    \]
    否则 $n(p,p')=0$。

    \item 对所有 $p,p'\in P$ 和 $g\in\Gamma$，有
    \[
      n(g\cdot p,g\cdot p')=n(p,p').
    \]

    \item 对每个\termFormalSum{}
    \[
      \beta=\sum_{q\in P}a_q q,
      \qquad a_q\in R,
    \]
    定义其\emph{支集}为
    \[
      \operatorname{supp}(\beta)
      =\{q\in P\mid a_q\ne0\}.
    \]
    现在定义\termFloerChainGroup{}
    \footnote{译者注：原书在本节三个\personNovikov{}有限性公式中都把条件排成
    “$\forall C\#(\cdots)<\infty$”，既没有写出 $C$ 的取值范围，也缺少分隔。
    译文依各式所用的实区间明示“对每个 $C\in\mathbb R$”；另外两处同样处理。}
    \[
      C_*(\mathfrak c)
      :=
      \left\{
        \beta=\sum_{q\in P}a_q q
        \ \middle|\
        \begin{gathered}
          a_q\in R,\quad
          \#\bigl(\operatorname{supp}(\beta)
            \cap\mathcal A^{-1}((C,\infty))\bigr)<\infty\\
          \text{对每个 }C\in\mathbb R
        \end{gathered}
      \right\}.
    \]
    于是，对每个 $p\in P$，形式和
    \[
      \partial p=\sum_{q\in P}n(p,q)q
    \]
    属于 $C_*(\mathfrak c)$。
    \footnote{译者注：原书在求和号下写成 $p\in P$，但被求和项及本句已经固定的
    输入分别使用 $q$ 与 $p$；译文把哑指标补正为 $q\in P$。}

    \item 定义 $\Gamma$ 的\termNovikovRing{}
    \footnote{译者注：原书定义的环元素写成 $\sum_{g\in\Gamma}b_g g$，有限性条件
    却取前一项链 $\beta$ 的 $\operatorname{supp}(\beta)$；译文将其补正为当前
    环元素 $\sigma$ 的支集。}为
    \[
      \Lambda^\uparrow_{\Gamma,\omega}
      :=
      \left\{
        \sigma=\sum_{g\in\Gamma}b_g g
        \ \middle|\
        \begin{gathered}
          b_g\in R,\quad
          \#\bigl(\operatorname{supp}(\sigma)
            \cap\omega^{-1}((-\infty,C))\bigr)<\infty\\
          \text{对每个 }C\in\mathbb R
        \end{gathered}
      \right\}.
    \]
    由 $\Gamma$ 在 $P$ 上的作用，$C_*(\mathfrak c)$ 以显然方式继承
    $\Lambda^\uparrow_{\Gamma,\omega}$-模结构。上面定义的算子
    \[
      \partial:P\longrightarrow C_*(\mathfrak c)
    \]
    延拓为一个 $\Lambda^\uparrow_{\Gamma,\omega}$-模同态
    \[
      \partial:C_*(\mathfrak c)\longrightarrow C_*(\mathfrak c),
    \]
    并且满足 $\partial\circ\partial=0$。
  \end{enumerate}
\end{enumerate}
\end{definition}

\stfhIndex{I027}{Zuoyongpu}{作用谱（action spectrum）}
\stfhIndex{I151}{Zhouqiqun}{周期群（period group）}
称 $\mathcal A$ 的像为复形 $\mathfrak c$ 的\emph{\termActionSpectrum{}}，称 $\omega$
的像为 $\Gamma$ 的\emph{\termPeriodGroup{}}；分别记作
\[
  \operatorname{Spec}(\mathfrak c)
  \quad\text{和}\quad
  \Gamma_\omega\subset\mathbb R.
\]

注意，\termNovikovFinitenessCondition{}蕴含每条\personFloer{}--\personNovikov{}链
\[
  \beta=\sum_{q\in P}a_q q,
  \qquad a_q\in R,
\]
都有一个“最高峰”，因而可以把它的\emph{层级}定义为
\[
  \ell(\beta)
  :=\max\{\mathcal A(q)\mid q\in\operatorname{supp}(\beta)\}.
\]
当 $\beta=0$ 时，约定 $\ell(\beta)=-\infty$。这定义了\emph{\termLevelFunction{}}
\footnote{译者注：原书先约定 $\ell(0)=-\infty$，紧接着却把 $\ell$ 的值域只写成
$\mathbb R$。译文补入该约定所需的扩充值 $-\infty$。}
\[
  \ell:C_*(\mathfrak c)\longrightarrow
  \mathbb R\cup\{-\infty\},
\]
并在 $C_*(\mathfrak c)$ 上诱导滤过
\[
  C_*^\lambda(\mathfrak c)
  :=\{\beta\mid \ell(\beta)\leq\lambda\}.
\]
由定义，映射 $\partial$ 保持该滤过，即
\[
  \partial\bigl(C_*^\lambda(\mathfrak c)\bigr)
  \subset C_*^\lambda(\mathfrak c).
\]
层级函数还在 $C_*(\mathfrak c)$ 上诱导一个拓扑；这个拓扑由如下\termNonArchimedeanMetric{}给出：
\begin{equation}
  d(\beta_1,\beta_2)
  =e^{\ell(\beta_1-\beta_2)}.
\label{eq:abstract-floer-nonarchimedean-metric}
\end{equation}

\stfhIndex{I191}{Fuzhi}{赋值（valuation）}
类似地，在\personNovikov{}环 $\Lambda^\uparrow_{\Gamma,\omega}$ 上定义\emph{\termValuation{}}
\[
  \nu(a)
  =\min\{\omega(g)\mid g\in\operatorname{supp}(a)\},
  \qquad
  a=\sum_{g\in\Gamma}b_g g\ne0,
\]
并约定 $\nu(0)=+\infty$
\footnote{译者注：原书没有定义 $\nu(0)$，但随后的度量在 $a_1=a_2$ 时需要
$e^{-\nu(0)}=0$。译文补出与该式相容的约定 $\nu(0)=+\infty$。}。它又在
$\Lambda^\uparrow_{\Gamma,\omega}$ 上诱导度量
\[
  d(a_1,a_2)=e^{-\nu(a_1-a_2)}.
\]
把以上讨论总结为如下命题。

\begin{proposition}
\label{prop:abstract-floer-continuity-equivariance}
映射
\[
  \partial:C_*(\mathfrak c)\longrightarrow C_*(\mathfrak c)
\]
关于\eqref{eq:abstract-floer-nonarchimedean-metric}中的距离 $d$ 连续。
$\Lambda^\uparrow_{\Gamma,\omega}$ 在 $C_*(\mathfrak c)$ 上的作用
\[
  \Lambda^\uparrow_{\Gamma,\omega}\times C_*(\mathfrak c)
  \longrightarrow C_*(\mathfrak c)
\]
也是连续的。此外，映射 $\partial$ 关于这个作用等变，因而该作用在
\[
  H_*(\mathfrak c,\partial)
\]
上诱导一个 $\Lambda^\uparrow_{\Gamma,\omega}$-模结构。
\end{proposition}

限制
\[
  \partial:C_*^\lambda(\mathfrak c)\longrightarrow
  C_*^\lambda(\mathfrak c)
\]
以及包含链映射
\[
  \iota:C_*^\lambda(\mathfrak c)\longrightarrow C_*(\mathfrak c)
\]
诱导映射
\[
  \iota_*:H_*^\lambda(\mathfrak c)\longrightarrow H_*(\mathfrak c).
\]
这使我们可以引入如下定义。

\begin{definition}
\label{def:abstract-spectral-number}
\stfhIndex{I176}{Pushu}{谱数（spectral number）}
对 $a\in H_*(\mathfrak c)$，定义 $a$ 的谱数为
\[
  \rho(a)
  =\inf\left\{
    \lambda\in\mathbb R
    \ \middle|\
    a\in\operatorname{Im}\bigl(
      \iota_*:H_*^\lambda(\mathfrak c)\to H_*(\mathfrak c)
    \bigr)
  \right\}.
\]
\end{definition}

等价地，可以把 $\rho(a)$ 定义为如下极小极大值：
\begin{equation}
  \rho(a)
  =\inf\{\ell(\beta)\mid
    \beta\in C_*(\mathfrak c),\ [\beta]=a\},
\label{eq:abstract-spectral-number-minimax}
\end{equation}
其中 $[\beta]$ 是 $\beta$ 的同调类。显然有 $\rho(0)=-\infty$。不过，即使
$a\ne0$，$\rho(a)>-\infty$ 是否成立也并非完全显然；而且当
$\rho(a)>-\infty$ 时，其值能否由一个代表同调类 $a$ 的闭链 $\beta$ 实现，也并非
完全显然。在经典临界点理论中，$\rho(a)$ 的有限性对应于代表给定同调类 $a$
的闭链所具有的某种“\termLinkingProperty{}” (BnR79)，而谱数 $\rho(a)$ 的可实现性
则对应于作用量泛函 $\mathcal A$ 的某种“\termPalaisSmaleCondition{}”。在这个一般
抽象框架中，\personUsher{}在 (Ush08) 中证明了这两个性质。

\begin{theorem}[\personUsher{} (Ush08)]
\label{thm:usher-abstract-spectrality}
对任意 $0\ne a\in H_*(\mathfrak c)$，有 $\rho(a)>-\infty$，并且存在闭链
$\alpha\in C_*(\mathfrak c)$，使得
\[
  \rho(a)=\ell(\alpha).
\]
特别地，对所有 $0\ne a\in H_*(\mathfrak c)$，都有
\footnote{译者注：原书在“特别地”之后把范围扩大为所有
$a\in H_*(\mathfrak c)$，但上一段已经给出 $\rho(0)=-\infty$，而
$-\infty$ 一般不属于作用谱。译文依定理首句保留条件 $a\ne0$。}
\[
  \rho(a)\in\operatorname{Spec}(\mathfrak c).
\]
\end{theorem}

这推广了 (Oh05c, Oh09a) 中在哈密顿\personFloer{}同调语境下以更几何的方法证明的结果。
我们将在第 21.3 节于这个几何语境中给出本定理的证明；一般抽象框架中的证明则请
\stfhIndex{I042}{Bianjieshendu}{边界深度（boundary depth）}
读者参见 (Ush08)。证明所用的一个重要成分是\termBoundaryDepth{}的概念；\personUsher{}在 (Ush11) 中把它正式引入\personFloer{}理论。（关于同一概念在链层\personFloer{}理论中的一些
用法，见 (FOOO09)、(Oh09a)。）

\begin{definition}
\label{def:boundary-depth}
设 $(C,\partial)$ 是一个 $\mathbb R$-\termFilteredChainComplex{}。定义其边界深度
$\beta(\partial)$
\footnote{译者注：原书把候选深度的范围写成 $\beta\in\mathbb R$；例如当
$\partial=0$ 时，这会使右端成为 $-\infty$，而不是边界深度的非负值。译文补为
$\beta\in[0,\infty)$。}为
\[
  \beta(\partial)
  =\inf\left\{
    \beta\in[0,\infty)
    \ \middle|\
    \text{对每个 }\lambda\in\mathbb R,\quad
    C^\lambda\cap\partial C
    \subset\partial\bigl(C^{\lambda+\beta}\bigr)
  \right\}.
\]
\end{definition}

下面证明一个由 $\Lambda^\uparrow_{\Gamma,\omega}$ 的有限性条件导出的有用结构
定理。首先有正合序列
\[
  0\longrightarrow\ker\omega
  \longrightarrow\Gamma
  \xrightarrow{\ \omega\ }\Gamma_\omega
  \longrightarrow0.
\]
因此，可以把
\[
  \sigma=\sum_{g\in\Gamma}a_g g
  \in\Lambda^\uparrow_{\Gamma,\omega}
\]
改写成
\[
  \sigma
  =\sum_{\lambda\in\Gamma_\omega}
    \left(
      \sum_{\substack{g\in\Gamma\\\omega(g)=\lambda}}a_g g
    \right).
\]
注意，每个内和
\[
  \sum_{\substack{g\in\Gamma\\\omega(g)=\lambda}}a_g g
\]
都是有限和。

\begin{proposition}
\label{prop:novikov-ring-period-group-coefficients}
考虑环
\[
  K=R[\ker\omega],
\]
并定义环
\[
  \Lambda^\uparrow_{\Gamma_\omega,K}
  :=
  \left\{
    \sum_{\lambda\in\Gamma_\omega}p_\lambda T^\lambda
    \ \middle|\
    \begin{gathered}
      p_\lambda\in K,\quad
      \#\{\lambda\in\Gamma_\omega\cap(-\infty,C)
        \mid p_\lambda\ne0\}<\infty\\
      \text{对每个 }C\in\mathbb R
    \end{gathered}
  \right\}.
\]
则映射
\[
  \sigma=\sum_{g\in\Gamma}a_g g
  \longmapsto
  \sum_{\lambda\in\Gamma_\omega}p_\lambda(\sigma)T^\lambda
\]
是一个环同构。
\footnote{译者注：原书没有定义 $p_\lambda(\sigma)$。按 $\omega$ 的纤维分组所得
的内和位于群环的相应陪集分量，并不直接属于
$K=R[\ker\omega]$。要得到所称的环同构，需要选取上述正合序列的一个群同态分裂
$s:\Gamma_\omega\to\Gamma$，再令
$p_\lambda(\sigma)=\sum_{\omega(g)=\lambda}
a_g\,s(\lambda)^{-1}g\in K$。这里 $\Gamma_\omega$ 是有限生成无挠阿贝尔群，故这种
分裂存在，但一般不典范；译文保留原书的公式并明示所缺选取。}
\end{proposition}
